\documentclass[a4paper,12pt,twoside,openright,titlepage]{book}

%Additional packages
\usepackage[utf8]{inputenc}
\usepackage[T1]{fontenc}
\usepackage[dutch,english]{babel}
\usepackage{imakeidx}
\usepackage{syntonly}
\usepackage[official]{eurosym}
%\usepackage[graphicx]
\usepackage{graphicx}
\graphicspath{ {./images/} }
\usepackage{float}
\usepackage{xurl}
\usepackage{hyperref}
\hypersetup{colorlinks=true, linkcolor=blue, citecolor=blue, filecolor=blue, urlcolor=blue, pdftitle=, pdfauthor=, pdfsubject=, pdfkeywords=}
\usepackage{tabularx}
\usepackage[table]{xcolor} % Table colors
\usepackage{scrextend}
\addtokomafont{labelinglabel}{\sffamily}
\usepackage{listings}
\usepackage{adjustbox}
\usepackage{color}
\usepackage{csquotes}

% Define colors
\definecolor{ashgrey}{rgb}{0.7, 0.75, 0.71}

% Listing style
\lstset{
  backgroundcolor=\color{ashgrey}, % choose the background color; you must add \usepackage{color} or \usepackage{xcolor}; should come as last argument
  basicstyle=\footnotesize,        % the size of the fonts that are used for the code
  breakatwhitespace=true,          % sets if automatic breaks should only happen at whitespace
  breaklines=true,                 % sets automatic line breaking
  extendedchars=true,              % lets you use non-ASCII characters; for 8-bits encodings only, does not work with UTF-8
  frame=single,	                   % adds a frame around the code
  rulecolor=\color{black},         % if not set, the frame-color may be changed on line-breaks within not-black text (e.g. comments (green here))
  keepspaces=true,                 % keeps spaces in text, useful for keeping indentation of code (possibly needs columns=flexible)
  columns=fullflexible,		   % make copy and paste possible
  showstringspaces=false,          % if true show spaces in strings adding particular underscores
  showspaces=false,                % if true show spaces everywhere adding particular underscores; it does not override 'showstringspaces'
}

% Uncomment for production
% \syntaxonly

% Style
\pagestyle{headings}

% Turn on indexing
\makeindex[intoc]

% Define document
\author{D. Leeuw}
\title{TCP: DNS}
\date{\today\\v.0.8.1}

\begin{document}
\selectlanguage{dutch}

\maketitle

\copyright\ 2023 Dennis Leeuw\\

\begin{figure}
\includegraphics[width=0.3\textwidth]{CC-BY-SA-NC.png}
\end{figure}

\bigskip

Dit werk is uitgegeven onder de Creative Commons BY-NC-SA Licentie en laat anderen toe het werk te kopi\"eren, distribueren, vertonen, op te voeren, en om afgeleid materiaal te maken, zolang de auteurs en uitgever worden vermeld als maker van het werk, het werk niet commercieel gebruikt wordt en afgeleide werken onder identieke voorwaarden worden verspreid.


%%%%%%%%%%%%%%%%%%%
%%% Introductie %%%
%%%%%%%%%%%%%%%%%%%

\frontmatter
\chapter{Over dit Document}
\section{Leerdoelen}
Dit document gaat over ethernet. Het leert je wat ethernet is. Om meer te weten te komen over wat er op je netwerk gebeurt maak je ook kennis met de tools: \texttt{arp} en \texttt{wireshark}.

Na het bestuderen van dit document heb je een basis kennis van ethernet, weet je wat het adres-formaat is van ethernet en kan je informatie over het netwerk opvragen via verschillende tools.

In de praktijk opdrachten leer je om je eigen kabel te maken en deze te testen. Tevens maak je je eigen Wi-Fi netwerk.

\section{Voorkennis}
Er wordt van de lezer verwacht dat deze kennis heeft van de volgende onderwerpen:
\begin{itemize}
\item TCP/IP
\item CMOS en BIOS tijd
\end{itemize}


%%%%%%%%%%%%%%%%%
%%% De inhoud %%%
%%%%%%%%%%%%%%%%%
\tableofcontents

\mainmatter
\chapter{Inleiding}
TCP/IP is het protocol van het Internet en tegenwoordig ook van de meeste andere computer-netwerken. Vroeger waren er veel verschillende protocollen. Van die verschillende protocolen is eigenlijk alleen TCP/IP overgebleven.

TCP/IP is ontworpen als een protocol dat moest werken over onbetrouwbare netwerkverbindingen en dat ook bij een kernwapen-oorlog waarbij delen van het netwerk uitvallen toch nog zijn data moet kunnen versturen van A naar B. Het moet dus feitelijk zelf uitzoeken hoe data het beste van A naar B kan komen.


\chapter{Het Domain Name System}
\section{De hosts-file}
De eerste methode die werd gebruikt om mensen de mogelijkheid te geven om namen voor computers te gebruiken en machines een IP-adres was de hosts-file. Het bestand bestaat nog steeds om lokaal machines van een naam te voorzien. Het hosts-bestand is een simpel tekst bestand dat er bijvoorbeeld zo uit kan zien:

\begin{lstlisting}[language=bash]
127.0.0.1	localhost
127.0.1.1	hcl01.localdomain	hcl01

# The following lines are desirable for IPv6 capable hosts
::1     localhost ip6-localhost ip6-loopback
ff02::1 ip6-allnodes
ff02::2 ip6-allrouters

10.181.65.42 my-test-host.example.com my-test-host
\end{lstlisting}

Je zou dit bestand kunnen zien als een kleine database met aan de linkerkant het IP-adres en rechts, gescheiden door wit (tab of spatie), de naam van de machine.

De hosts-file kan je op Linux of Mac OS X systemen terug vinden als \texttt{/etc/hosts} op een Windows systeem is het \texttt{C:\textbackslash Windows\textbackslash System32\textbackslash Drivers\textbackslash etc\textbackslash hosts}. Het feit dat op een Windows systeem ook het laatste stukje \texttt{\textbackslash etc\textbackslash hosts} is heeft te maken met het feit dat Microsoft de BSD TCP/IP-stack heeft gebruikt als basis voor de TCP/IP-drivers in Windows.

Als je de laatste regel van het voorbeeld (wat begint met 10.181.65.42) toevoegt aan je eigen hosts file dan kan je met
\begin{lstlisting}[language=bash]
ping my-test-host
\end{lstlisting}
zien dat de machine contact probeert te maken met 10.181.65.42. Het zal dit waarschijnlijk niet kunnen doen, maar je ziet dat hij de naam wel omzet naar het opgegeven IP-adres.

Voordat je verder gaat met dit document verwijder je eerst de toegevoegde regel weer uit je hosts-file.


\section{DNS}
Op een lokaal netwerk is het mogelijk om het hosts-bestand bij te houden, maar bij een wereldomvattend netwerk zoals het Internet is dat niet meer te doen. Er moest dus een wereldwijde database komen die op afstand geraadpleegd kan worden. Het nadeel van een enkele database is dat het tevens een single point of failure (SPOF) is. Als de database-server omvalt is er geen mogelijkheid meer om namen om te zetten naar IP-adressen. Ook en verbroken netwerkverbinding kan zorgen dat een deel van het netwerk niet meer bij de database kan. Om al deze problemen het hoofd te kunnen bieden is er gekozen voor een gedistribureerd systeem. Er is niet \'e\'en enkele database dat de namen van de machines kent.

Zoals je misschien wel weet is net Internet opgedeeld in zogenaamde domeinen. Het adres www.google.com bestaat uit drie delen. Het eerste deel is www, dit is de host-naam. Het google deel is het domein en het com deel is het top-level-domain (TLD). De database-servers worden dan ook Domain Name Servers genoemd, afgekort DNS.

Er zijn verschillende TLD's op de wereld. Elk land heeft zijn eigen TLD, die voor Nederland is nl, voor Duitsland de en die voor Belgi\"e be. Zo heeft elk land zijn eigen TLD. Daarnaast zijn er vele generieke TLD's zoals com voor commerci\"ele bedrijven, org voor organisaties en edu voor educatieve instellingen. Elke TLD heeft zijn eigen name-server wat al zorgt voor een enorme wereldwijde spreiding.

Een organisatie als Google heeft zijn eigen domein (google.com) en daarmee ook zijn eigen domein-database (DNS). Dus als je wilt weten wat het adres is dat hoort bij www.google.com dan vraag je dat aan de DNS van Google.

%%%% Van onder naar boven uitzoeken wie het is

%%%% Hoe weet een client welke DNS server hij moet hebben (geen eigen dns-server)

\section{De werking van DNS}
Hoe zorgt een client, bijvoorbeeld je browser, dat je bij de domeinnaam een IP adres krijgt?

Je browser geeft aan de TCP/IP stack door dat hij bijvoorbeeld verbinding wil maken met www.microsoft.com. Het eerste dat je TCP/IP stack dan gaat doen is kijken of het weet welke name servers er zijn. Deze name server moet je handmatig hebben geconfigureerd of ze moeten door DHCP aan je machine zijn gegeven.

Je machine neemt contact op met een DNS op port 53 en vraagt aan de server of hij weet wat het IP-adres is van www.microsoft.com.

De kans is natuurlijk klein dat jouw name server de IP-adressen van Microsoft kent. Dus je name server neemt contact op met een root name server. Een root name server kent de Top Level Domains (TLDs) van het Internet. De root server kan onze server verwijzen naar de name server voor het .com domain. Die name server kan ons verwijzen naar het IP-adres van de name server van het .microsoft.com domein.

Als we het IP-adres hebben van de name server van Microsoft dan kan onze name server aan de name server van Microsoft vragen wat het IP adres is voor www.microsoft.com en dat antwoord kan onze server dan weer aan onze machine gegeven. Daarna kunnen we een verbinding leggen met de web-server van Microsoft.

Om toekomstige DNS queries te versnellen, kunnen de gevonden oplossing tijdelijk opslaan (cache) zodat we een volgende query sneller kunnen beantwoorden. Voor veel gebruikte domeinen helpt dit om het proces te versnellen. Het nadeel is natuurlijk wel dat als er in de DNS bij Microsoft een wijziging plaats vindt wij die niet meteen medelen aan onze gebruikers. Pas als de cache rijd verlopen is zullen we bij Microsoft controleren of het gevonden IP nog geldig is en tot de conclusie komen dat dit niet het geval is. Pas daarna weten ook onze gebruikers het.


\section{Root Hints}
De InterNIC is verantwoordelijk voor het bijhouden van de root van de name servers van het Internet. Dit is een kort lijstje van alle servers wereldwijd die alle TLDs kennen. Dit is een lijstje dat je af en toe eens moet updaten om te zorgen dat je de laatste versie hebt. Het is beschikbaar op \url{https://www.internic.net/domain/named.cache}. Het is bekend onder de naam \textquote{hints-file}.

Het bestandje heeft allemaal regels die vertellen waar root servers gevonden kunnen worden. De root van een domein is de . (punt). Eigenlijk zou je een domein naam dus moeten schrijven als wwww.microsoft.com. met een punt aan het einde. Maar omdat die punt er altijd moet staan laten we hem weg en nemen we aan dat die er is. In de hints-file komen we deze punt weer tegen:
\begin{lstlisting}[language=bash]
.                        3600000      NS    A.ROOT-SERVERS.NET.
A.ROOT-SERVERS.NET.      3600000      A     198.41.0.4
A.ROOT-SERVERS.NET.      3600000      AAAA  2001:503:ba3e::2:30
\end{lstlisting}
We zien hier dat voor de root (.) we de name server (NS record) A.ROOT-SERVERS.NET. moeten gebruiken en dat die server het IPv4 adres (A record) 198.41.0.4 heeft. Ook een IPv6 adres (AAAA record) is gegeven voor de name server.



\chapter{Forwarding DNS}
De meest simpele vorm van DNS is het forwarden van een request aan een volgende DNS. Dit gebeurd vaak in de wat kleinere netwerken. In interne DNS server wordt door alle clients gebruikt om hun requests naar toe te sturen en de interne DNS forward de request naar de DNS van de provider. De interne DNS werkt dan alleen als cache zodat veel voorkomende queries niet allemaal bij de DNS van de provider terecht komen. Dus hoewel het toevoegen van een extra DNS een extra stap is en dus vertraging oplevert kan het door de caching uiteindelijk toch een verbetering zijn.



\chapter{De DNS database}
\section{Serial}
De database van de verschillende DNSen kan wijzigen. Het kan bijvoorbeeld gebeuren dat de web-server op een andere machine komt te draaien en dat dus het IP-adres wijzigt. Dit willen we in het name-server database wijzigen, zolang de andere name-servers aan ons systeem vragen wat het IP adres is gaat het goed, dan is de wijziging onmiddelijk, maar als een name-server een cache heeft aangelegd, dan zal die voorlopig zijn cache blijven gebruiken. Het kan dus enige tijd duren voordat een wijziging op het gehele Internet bekend is.

De manier om caches te laten verlopen is het updaten van de \textquote{serial}. De serial is een numeriek ID. Een hogere serial betekent een nieuwere versie. Voor de serial wordt meestal de datum + een versie nummer genomen: 2023120603. Deze serial staat voor 6 december 2023 de derde wijziging. Je kan zo maximaal 99 wijzigingen op een dag doen (01 tot 99).

Een DNS kan dus de serial opvragen om te weten te komen of het opnieuw een query moet doen naar een IP-adres of dat hij de huidige cache kan laten bestaan.


\section{NS record}
Het name server (NS) record vertelt wat de name-servers zijn voor een domein. Let op dat je bij dit record niet 127.0.0.1 invult, maar altijd het externe IP-adres van je server, daar waar iedereen bij kan. Als je je eigen Internet domein hebt, moet je name server ook werkelijk aan het Internet hangen.


\section{MX record}
De wereld wil ook graag weten hoe e-mail naar een organisatie verstuurd kan worden. Hiervoor is er een speciaal veld in de DNS database, het zogenaam de MX ofwel Mail eXchange record. Het MX-record verwijst naar de hostname (plus domainname) van de mail server en geeft tevens een prioriteit mee. De mail-server met de laagste prioriteit moet gebruikt worden, de andere mag alleen gebruikt worden als de eerste niet bereikbaar is of als deze aangeeft dat hij niet gebruikt mag worden.


\section{A record}
Het A record in de DNS database staat voor een adres-record. In dit veld wordt de hostname opgegeven met het bijbehorende IP-adres.


\section{PTR record}
Het kan voorkomen dat we ook bij een IP-adres willen weten welk domein erbij hoort. Stel bijvoorbeeld dat iemand op ons systeem inlogt vanaf 142.251.36.46. Om dat mogelijk te maken met de tree-werking van DNS moeten we ook een reverse tree hebben. De oplossing die gekozen is om een apart domein te maken voor reverse lookups. Het domein is in-addr.arpa. Het TLD is dus .arpa. en de daarbij behorende functionele domein is in-addr (Internet Address).

Domeinen zoeken we op van TLD naar domein naar host. Om IP-adressen te resolven moeten we dat ook omgekeerd doen, dus voor 142.251.36.46 moet dat worden eerst 142, dan 251, vervolgens 36 om uit te komen bij 42. Dus de query zou moeten zijn 46.36.251.142.in-addr.arpa. Dat is een moeilijke om elke keer goed te doen, dus lossen de tools die we gebruiken om DNS-queries uit te voeren dit zelf op. Ook de DNS database lost dit zelf op, maar we moeten hem wel vertellen dat het een reverse-database is.

De reverse database is dus een andere database dan de forward database. In de reverse database hebben we dan ook geen A-records, maar PTR records. Het Pointer-record zegt hij dus behoort tot het in-addr.arpa domain.



\chapter{Opdrachten}
\section{nslookup}
Een tool die je op bijna elk operating system kan gebruiken om vragen te stellen aan DNS is \texttt{nslookup}. In de meest eenvoudige vorm kan \texttt{nslookup} gebruikt worden om een IP-adres bij een domeinnaam op te vragen:
\begin{lstlisting}[language=bash]
nslookup www.google.com
\end{lstlisting}

Je kan ook specifiek om het IPv4 adres vragen door alleen het A-record op te vragen:
\begin{lstlisting}[language=bash]
nslookup -type=A www.google.com
\end{lstlisting}

We weten dat DNS aan de root van de DNS-tree begint en langzaam verder gaat. Dat gaan wij nu ook doen:
\begin{enumerate}
\item Gebruik \texttt{nslookup} om een met lijst name-servers op te vragen die behoren bij het com domain.
\item Vraag aan \'e\'en van de gevonden name-servers wat de name-servers zijn voor google.com.
\item Vraag aan \'e\'en van de name-servers van Google wat de IP-adressen zijn van de mail-servers van Google (2 queries).
\end{enumerate}

% dig
\section{DNS installatie op Windows}
Voor het opzetten van een DNS op Windows is een document van Microsoft dat je kan vinden op \url{https://learn.microsoft.com/en-us/windows-server/networking/dns/quickstart-install-configure-dns-server?tabs=powershell}. De installatie kan je doen op een fysiek systeem of op een virtual machine.


%%%%%%%%%%%%%%%%%%%%%
%%% Index and End %%%
%%%%%%%%%%%%%%%%%%%%%
\backmatter
\printindex
\end{document}

%%% Last line %%%
